\chapter{Metodologia}\label{chp:AP_Results}

%Use diagramas de arquitetura (ex.: fluxos de dados, containers Docker).

%Inclua tabelas de requisitos e exemplos de telas para ilustrar a metodologia.

%Sempre explique por que escolheu cada tecnologia, mostrando alinhamento com os objetivos do sistema.

\section{Tipo de Pesquisa}\label{sec:tipoPesquisa}
%Explique se o trabalho é aplicado, experimental ou de desenvolvimento tecnológico.

%Exemplo: pesquisa aplicada com desenvolvimento de protótipo de sistema para reserva de ambientes acadêmicos.


\section{Levantamento de Requisitos}\label{sec:requisitos}
%Descreva como você identificou as necessidades da comunidade acadêmica.
%Métodos: entrevistas, questionários, observação, análise de sistemas existentes.
%Liste os requisitos funcionais e não funcionais do sistema.
Para identificar as necessidades e expectativas dos usuários do sistema de reservas, foi elaborado um questionário direcionado aos funcionários da biblioteca, responsáveis pela gestão dos espaços acadêmicos. O objetivo é coletar informações sobre o uso atual das salas, mesas e auditórios, dificuldades enfrentadas no processo de reserva e funcionalidades desejadas para o novo sistema. O questionário completo encontra-se disponível no \hyperref[apendice:entrevistas]{Apêndice A}.

O questionário combina perguntas objetivas e abertas, permitindo tanto a obtenção de dados quantitativos, como frequência de uso e problemas mais comuns, quanto dados qualitativos, como sugestões de melhorias e funcionalidades essenciais. Entre os tópicos abordados, destacam-se:
\begin{itemize}
  \item Uso atual dos espaços: frequência de reservas, problemas encontrados, sistemas atualmente utilizados.
  \item Funcionalidades desejadas: prioridades na gestão de reservas, tipos de visualização dos ambientes e recursos que devem ser disponibilizados.
  \item Usabilidade e experiência do usuário: facilidade de uso, informações essenciais na interface e percepção sobre a redução de conflitos.
  \item Sugestões e melhorias: opiniões abertas sobre funcionalidades essenciais e formas de otimizar o sistema.
\end{itemize}

A aplicação do questionário tem o potencial de identificar requisitos funcionais e não funcionais do sistema, bem como alinhar o desenvolvimento da interface e das funcionalidades às reais necessidades da comunidade acadêmica. Os dados coletados serão analisados e utilizados como base para o projeto do sistema, garantindo que a solução proposta seja eficiente, prática e adequada ao contexto da universidade.

\section{Arquitetura do Sistema}\label{sec:arquitetura}
%Apresente diagrama de arquitetura (Front-end, Back-end, Banco de dados).
%Explique a escolha de tecnologias: Angular (front-end), PHP (back-end), MySQL (banco de dados), Docker (containerização).
%Mostre como cada componente se comunica (ex.: Angular → PHP → MySQL).

\section{Desenvolvimento do Sistema}\label{sec:desenvolvimento}
\subsection{Front-end (Angular)}\label{subsec:angular}
%Componentes, rotas, interface de visualização das plantas baixas.
%Interação com o usuário (reserva de mesas, salas, auditorios).

\subsection{Back-end (.NET)}\label{subsec:dotnet}
%APIs e endpoints para gerenciamento de reservas.
%Validação de dados, autenticação de usuários.

\subsection{Banco de Dados (MySQL)}\label{subsec:mysql}
%Estrutura de tabelas, relacionamentos, integridade de dados.

\subsection{Docker}\label{subsec:docker}
%Containerização do ambiente de desenvolvimento e produção.
%Benefícios de portabilidade e padronização.


\section{Interface de Visualização dos Espaços}\label{sec:interface}
%Explique como você representou as plantas baixas de forma fiel.
%Métodos de renderização (2D/3D, grid, mapas interativos).
%Integração entre front-end e dados do banco de dados.

\subsection{Sub-tópico}\label{subsec:landsat}

\section{Testes e Validação}\label{sec:testesValidacao}
%Testes funcionais: reserva de mesas, salas e auditorios.
%Testes de usabilidade: interface intuitiva e responsiva.
%Testes de integração: front-end, back-end e banco funcionando corretamente.

\subsection{Sub-tópico}\label{subsec:landsat}

\section{Ferramentas e Ambientes de Desenvolvimento}\label{sec:ferramentas}
%Softwares utilizados (VS Code, MySQL Workbench, Docker Desktop, Git, etc.).

%Ambientes (desenvolvimento local, testes, produção).

\pagebreak

